\section{Model Results and Clinical Interpretation}

The joint longitudinal--interval model was fitted to 87 patients with 495 longitudinal MRD observations and 275 at-risk DMR intervals. Sixty-eight patients reached DMR and 19 were censored. The data structure reflected real-world CML monitoring: median follow-up was 39.0 months, the median visit gap was 6.0 months, 50.3\% of visit gaps exceeded 6 months, and 239 longitudinal observations were at or below the log-MRD assay floor. These features support a model that jointly handles repeated molecular measurements, irregular monitoring, assay-floor censoring, and interval-observed DMR onset.

The Bayesian computation was adequate for substantive interpretation. The fitted model used four chains with 1000 warmup and 1000 sampling iterations per chain, yielding 4000 post-warmup draws. There were no divergent transitions, the maximum treedepth was 7, and the minimum approximate E-BFMI was 0.500. The main fixed-effect and event-association parameters had acceptable convergence diagnostics. The random-effect correlation was weakly estimated (\(\Omega_{12}=0.026\), 95\% posterior interval -0.417 to 0.621; R-hat 1.019; bulk ESS 191), and was therefore treated as a nuisance parameter rather than a substantive clinical finding.

The longitudinal submodel showed a marked decline in latent log-MRD over follow-up. The posterior mean for the log-time effect was -3.539 (95\% posterior interval -4.497 to -2.613), indicating strong molecular decline over time. The quadratic log-time term was positive on average (posterior mean 0.499, 95\% posterior interval -0.005 to 1.012), suggesting nonlinear response dynamics, although the 95\% interval nearly included zero. The bone marrow sample-source effect was positive on average but uncertain (posterior mean 0.613, 95\% posterior interval -0.826 to 2.065), and should be interpreted as measurement-source adjustment rather than a confirmed clinical effect.

The interval DMR submodel provided the main clinical signal. The association parameter linking latent log-MRD to DMR was negative, with posterior mean -1.256 and 95\% posterior interval -1.851 to -0.809. Under the fitted parameterization, a one-unit lower latent log-MRD value corresponded to an approximately 3.5-fold higher interval DMR hazard component. Thus, lower estimated molecular burden was strongly associated with higher probability of reaching DMR, supporting the biological and clinical coherence of the joint model.

Posterior mean interval-level DMR probabilities showed substantial heterogeneity across at-risk intervals. The median interval probability was 0.136, with interquartile range 0.039 to 0.383; 51 intervals had posterior mean probability below 0.01, whereas 20 intervals had posterior mean probability above 0.80. This separation supports the potential clinical value of the model as a dynamic monitoring and risk-stratification framework. However, these probabilities should not yet be used as treatment-decision thresholds without external validation and calibration assessment.

Overall, the model contributes a clinically oriented statistical framework for CML molecular monitoring. It uses the full serial MRD history, respects irregular visit timing, accounts for assay-floor measurements, and links evolving molecular burden to interval-observed DMR. The results support individualized monitoring and future dynamic prediction, but the current single-cohort analysis should be framed as a foundation for clinical decision support rather than a validated decision rule.
